\documentclass[11pt]{article}
\usepackage[utf8]{inputenc}
\usepackage{amsmath}
\usepackage{amssymb}
\usepackage{geometry}
\geometry{margin=1in}
\usepackage{enumitem}
\setlength{\parindent}{0pt}
\setlength{\parskip}{6pt}

\begin{document}

\begin{center}
\Large\textbf{New Practice Exam}
\end{center}

\vspace{0.5cm}

\textbf{Instructions:}
\begin{itemize}
\item Read each question carefully.
\item Show all your work for full credit.
\item You may use a calculator unless otherwise specified.
\end{itemize}

\vspace{0.5cm}

\begin{enumerate}[label=\textbf{Q\arabic*.}]
    \item [\textbackslash\{\}textbf\{Lucky Dice\}, 10 points] John loves playing games with dice. He has a special 12-sided die, with faces numbered from 1 to 12. Each time John rolls the die, the outcome is uniformly random, independent of other rolls.
 \textbackslash\{\}begin\{enumerate\}[label=(\textbackslash\{\}alph*)]
 \textbackslash\{\}item What is the expected value of the number he rolls?
 \textbackslash\{\}item If John rolls the die 24 times, what is the expected number of times he rolls a number greater than 8?
 \textbackslash\{\}item John rolls the die until he gets a 12. What is the expected number of rolls he makes?
 \textbackslash\{\}item John rolls the die 36 times. What is the expected number of unique numbers he gets?
 \textbackslash\{\}item Suppose John wants to roll a 7. Use the Most Used Inequality, \$1-x\textbackslash\{\}leq e\textasciicircum{}\{-x\}\$, to show that if he rolls the die 12 times, the probability that he gets a 7 at least once is at least \$1/2\$.
 \textbackslash\{\}end\{enumerate\}

    \item [\textbackslash\{\}textbf\{Random Shopping\}, 10 points] After a long week of work, Sarah decides to go shopping. She has a list of 10 items to buy. Each item has a 0.5 chance of being in stock, independently of other items.
 \textbackslash\{\}begin\{enumerate\}[label=(\textbackslash\{\}alph*)]
 \textbackslash\{\}item What is the expected number of items Sarah will be able to buy?
 \textbackslash\{\}item What is the probability that Sarah will be able to buy all the items on her list?
 \textbackslash\{\}item If each item costs \textbackslash\{\}\$20, what is the variance of the total amount Sarah will spend?
 \textbackslash\{\}item Sarah will buy herself a treat if she can buy at least 8 items from her list. What is the expected value of the indicator that Sarah gets a treat?
 \textbackslash\{\}end\{enumerate\}

    \item [\textbackslash\{\}textbf\{Random Walk\}, 10 points] Alice decides to take a random walk. She starts at point 0 and with each step, she either moves one step to the right with probability 0.6 or one step to the left with probability 0.4, independently of other steps.
 \textbackslash\{\}begin\{enumerate\}[label=(\textbackslash\{\}alph*)]
 \textbackslash\{\}item What is the expected value of her position after 10 steps?
 \textbackslash\{\}item What is the variance of her position after 10 steps?
 \textbackslash\{\}item What is the probability that Alice returns to the starting point after 10 steps?
 \textbackslash\{\}end\{enumerate\}

    \item [\textbackslash\{\}textbf\{News Errors\}, 10 points] Michael is a journalist who writes 5 articles every day; for each article he writes, he makes a mistake with probability 0.1 independently of other articles. Rachel is a trainee journalist who writes 3 articles every day; for each article she writes, she makes a mistake with probability 0.2 independently of other articles.
\textbackslash\{\}begin\{enumerate\}[label=(\textbackslash\{\}alph*)]
 \textbackslash\{\}item Find the standard deviation of the total number of mistakes they make in 7 days.
 \textbackslash\{\}item The two journalists make a total of 3 mistakes on a given day. What is the probability that Michael made more mistakes than Rachel that day?
 \textbackslash\{\}item The next day, at least one of the 8 articles they wrote contains a mistake. Find the expected number of total mistakes they made that day.
\textbackslash\{\}end\{enumerate\}

    \item [\textbackslash\{\}textbf\{Lottery Tickets\}, 10 points] A lottery has 100 tickets, numbered from 1 to 100. A ticket is randomly picked as the winning ticket. Let \$X\$ be the number on the winning ticket.
 \textbackslash\{\}begin\{enumerate\} [label=(\textbackslash\{\}alph*)]
 \textbackslash\{\}item Define the distribution of \$X\$. What is the expected value of the number on the winning ticket?
 \textbackslash\{\}item What is the probability that the number on the winning ticket is between 50 and 70?
 \textbackslash\{\}item Let \$Y\$ be the indicator random variable that takes the value 1 if the winning ticket has a number between 50 and 70, and 0 otherwise. Find the distribution, parameter, and PDF of \$Y\$.
 \textbackslash\{\}item Suppose John, Sarah, Michael, and Rachel each randomly pick a ticket from the lottery. What is the expected number of times a ticket numbered between 50 and 70 is picked?
 \textbackslash\{\}end\{enumerate\}

\end{enumerate}

\end{document}